

\newpage

%\footline={\hss\folio\hss}
%\pageno=1

\dspace
\subhead
Appendice. Rappels sur les groupes
alg\'ebriques unipotents
\endsubhead

Dans tout cet appendice, nous supposons
\^{e}tre en caract\'eristique z\'ero.

\subhead
A.1
\endsubhead
En caract\'eristique z\'ero, le th\'eor\`eme de
Poincar\'e-Birkhoff-Witt peut \^{e}tre
\'enonc\'e comme suit. 
Pour $\scrL$ une
alg\`ebre de Lie sur $K$, soit
$\varphi_n\colon\,\Sym^n\scrL\to\scrU\scrL$ le
produit sym\'etris\'e, de la puissance
symm\'etrique $n^{\ieme}$ de $\scrL$ dans
l'alg\`ebre enveloppante de $\scrL$.
Pour $x$ dans $\scrL$, il envoie
$x^n\in\Sym^n\scrL$ sur $x^n\in\scrU\scrL$.
Le morphisme d'espaces vectoriels
$$
\varphi_*\colon\,\Sym^*\scrL\to\scrU\,\scrL
\tag{A.1.1}
$$
est un isomorphisme de cog\`ebres.
Il transforme la filtration croissante de
$\Sym^*\scrL$ par les $\oplusop_{n\Le
p}\Sym^n\scrL$ en la filtration de
$\scrU\scrL$ par les images des
$\oplusop_{n\Le p}\otimes^n\scrL$.

Sous cette forme, le th\'eor\`eme est valable
dans toute cat\'egorie tensorielle (de
caract\'eristique $0$) pseudo-ab\'elienne
(pour donner un sens \`a $\Sym^n$) qui admet
des sommes d\'enombrables. 
Le cas qui nous int\'eresse est celui des
espaces vectoriels filtr\'es.
Soit $\scrL$ une alg\`ebre de Lie munie d'une
filtration d\'ecroissante $F$.
La filtration correspondante de $\scrU\,\scrL$
est caract\'eris\'ee par la propri\'et\'e qu'une
filtration sur un $\scrL$-module $M$ est
compatible \`a la filtration de $\scrL$ si et
seulement si elle l'est \`a celle de
$\scrU\,\scrL$.
La version cat\'egorique de
Poincar\'e-Birkhoff-Witt implique
 la compatibilit\'e de (A.1.1) aux filtrations.

Une filtration $F$ sur un
$\scrL$-module $M$ est compatible la
filtration centrale descendante $\grZ$ de
$\scrL$, si et
seulement si les $F^pM$ sont des sous-modules
et que $\scrL$ agit trivialement sur les
$\Gr_F^pM$.
La filtration de $\scrU\,\scrL$ induite par le
filtration $\grZ$ de $\scrL$ est donc la
filtration par les puissances de l'id\'eal
d'augmentation $I$.
Par compatibilit\'e de (A.1.1) aux
filtrations, on a

\proclaim{Corollaire A.2}
Si l'alg\`ebre de Lie $\scrL$, de s\'erie
centrale descendante  $\grZ$, est nilpotente
de classe $c$ $(\grZ^{c+1}=0)$, alors, si on
identifie $\Sym^*\scrL$ et $\scrU\,\scrL$ par
(A.1.1), on a
$$
I^n\Supset\oplusop_{p\Ge n}\Sym^p\scrL\Supset
I^{nc}\tag{A.2.1}
$$
\endproclaim

Notons $\scrU^{\uphat}\scrL$ le compl\'et\'e
$I$-adique de $\scrU\,\scrL$. 

\proclaim{Corollaire A.3}
Si l'alg\`ebre de Lie $\scrL$ est nilpotente,
l'isomorphisme (A.1.1) induit par compl\'etion
un isomorphisme d'espace vectoriels
$$
\prod\limits_{n\Ge 0}\Sym^n\scrL\to
\scrU^{\uphat}\scrL.
\tag{A.3.1}
$$
\endproclaim

L'isomorphisme (A.3.1) est compatible aux
topologies des deux membres.
A gauche, la topologie est la topologie
produit (d'espaces discrets).
A droite, celle de limite projective des
(espaces discrets) $\scrU\,\scrL/I^n$.

\subhead
A.4
\endsubhead
Le foncteur ``alg\`ebre de Lie'' est une
\'equivalence de la cat\'egorie des groupes
alg\'ebriques unipotents avec celle des
alg\`ebres de Lie nilpotentes de dimension
finie.
De plus, le foncteur $\rho\mapsto d\rho$ est
une \'equivalence de la cat\'egorie des
repr\'esentations lin\'eaires (de dimension
finie) de $V$, suppos\'e unipotent, avec celle
des repr\'esentations nilpotentes de son
alg\`ebre de Lie $\bfv$.

On dispose d'une application exponentielle
$$
\exp\colon\,\bfv\to V\tag{A.4.1}
$$
caract\'eris\'ee par la propri\'et\'e que pour
toute repr\'esentation lin\'eaire  $\rho$ de
$V$, on ait
$$
\rho(\exp(x))=\exp(d\rho(x)).\tag{A.4.2}
$$
C'est un isomorphisme de $\bfv$, muni de la
loi de groupe de Campbell-Hausdorff, avec $V$.

\subhead
A.5
\endsubhead
Soit $V$ un groupe alg\'ebrique unipotent
d'alg\`ebre de Lie $\bfv$.
L'alg\`ebre enveloppante $\scrU\bfv$ est
l'alg\`ebre des op\'erateurs diff\'erentiels
sur $V$ invariants par translations \`a gauche.
Comme cog\`ebre, c'est le dual topologique du
compl\'et\'e de l'anneau local de $V$ en $1$.
L'application exponentielle (A.4.1) induit un
isomorphisme de la cog\`ebre $\scrU\bfv$ avec
la cog\`ebre $\Sym^*\bfv$ des op\'erateurs
diff\'erentiels en $0$ sur l'espace affine
$\bfv$.

Par r\'eduction au cas o\`u $\bfv$ est de
dimension un, on v\'erifie que cet
isomorphisme envoie $x^n\in\Sym^n\bfv$ sur
$x^n\in\scrU\bfv$.
C'est donc l'isomorphisme (A.1.1).
Pour l'espace affine $\bfv$, l'alg\`ebre
affine de $\bfv$ est le dual topologique de la
cog\`ebre compl\'et\'ee
$\prod\limits_{n>0}\Sym^n\bfv$.
Appliquant A3, on obtient:

\proclaim{Proposition A.6}
L'alg\`ebre affine de $V$ est le dual
topologique de la cog\`ebre $\scrU\bfv$,
compl\'et\'ee pour la topologie $I$-adique.
\endproclaim

Dualement, le dual (muni de la topologie
faible) de l'alg\`ebre affine de $V$ est
$\scrU^{\uphat} \bfv$.
Les points de $V$ \`a valeurs dans une
alg\`ebre $A$ (homomorphismes
de l'alg\`ebre affine dans $A$) sont les {\it
\'el\'ements groupaux} $g$ de
$A\widehat{\otimes}\scrU^{\uphat}\bfv$: les
\'el\'ements d'augmentation $1$ tels que
$\Delta g=g\otimes g$.
Compatibilit\'es: l'application exponentielle
(A.4.1) est l'exponentielle, de $A\otimes g$
dans $V(A)\Subset
A\widehat{\otimes}\scrU^{\uphat}\bfv$.
La loi de groupe est le produit dans $\scrU^{\uphat}
\bfv$.

\subhead
A.7
\endsubhead
Un sch\'ema en groupe affine $G$ est la
limite projective de ses quotients $G_\alpha$
qui sont des groupes alg\'ebriques.
Il y a lieu de d\'efinir son alg\`ebre de Lie
$\Lie(G)$ comme \'etant la pro-alg\`ebre de
Lie limite projective des $\Lie(G_\alpha)$.
De fa\c{c}on \'equivalente, c'est la limite
projective des $\Lie(G_\alpha)$, qui sont de
dimension finie, munie de sa topologie de
limite projective d'espaces discrets.
C'est le dual de la limite inductive
$\colim(\Lie(G_\alpha)\upvee)$, muni de sa
topologie faible.

\subhead
A.8
\endsubhead
Un sch\'ema en groupe affine $G$ est {\it
pro-unipotent} s'il est limite projective de
groupes alg\'ebriques unipotents, i.e. si les
$G_\alpha$ de A.7 sont unipotents.
On d\'efinit l'alg\`ebre enveloppante
compl\'et\'ee $\scrU^{\uphat} \Lie(G)$ de
$\Lie(G)$ par
$$
\scrU^{\uphat}\Lie\,G=\lim\nolimits_\alpha
\scrU^{\uphat}\Lie(G_\alpha)=\lim\nolimits_{\alpha,n}
\scrU(\Lie(G_\alpha))/I^n,
$$
cette limite \'etant munie de sa topologie de
limite projective.
Avec ces d\'efinitions, A.6 reste valable,
comme on le v\'erifie par passage aux limites.

\subhead
A.9
\endsubhead
Cas particulier. Soit $\Lib((x_i)_{i\in I})$
l'alg\`ebre de Lie libre sur $K$ engendr\'ee
par une famille finie de g\'en\'erateurs
$x_i$.
Notons-la simplement $\Lib$.
Son alg\`ebre enveloppante est l'alg\`ebre
associative libre $K\left<(x_i)_{i\in I}\right>$.

Soit $\Lib^{\uphat}$ le compl\'et\'e de Lib pour
la filtration centrale descendante.
C'est l'alg\`ebre de Lie du sch\'ema en groupe
pro-unipotent
$$
\exp(\Lib^{\uphat})=\lim\,\exp(\Lib/\grZ^n).
$$
Son alg\`ebre enveloppante compl\'et\'ee est
l'alg\`ebre des s\'eries formelles
associatives $K\left<\!\left<(x_i)\right>\!\right>$.
Les points de $\exp(\Lib^{\uphat})$ \`a valeurs
dans une $K$-alg\`ebre $A$ sont les \'el\'ements
groupaux $g$ de 
$K\left<\!\left<(x_i)\right>\!\right>$:
terme constant $1$ et $\Delta g=g\otimes g$,
$\Delta$ \'etant le coproduit $K\left<\!\left<
(x_i)\right>\!\right>\to K\left<\!\left<
(x_i)\right>\!\right>\widehat{\otimes}K\left<\!\left<
(x_i)\right>\!\right>$ tel que $\Delta
x_i=x_i\otimes 1+1\otimes x_i$.

\subhead
A.10
\endsubhead
Soit $\bfv$ l'alg\`ebre de Lie d'un groupe
alg\'ebrique nilpotent $V$, et soit
$D$ une filtration
d\'ecroissante finie de l'alg\`ebre de Lie
$\bfv$, telle que $D^0(\bfv)=\bfv$.
On notera encore $D$ la filtration de $V$ par
les sous-groupes d'alg\`ebres de Lie les
$D^p\bfv$.
Puisque $[D^p\bfv, D^q\bfv]\Subset
D^{p+q}\bfv$, ce sont des sous-groupes
distingu\'es, et pour $p\Ge 1$, $D^pV/D^{p+1}V$
est ab\'elien.

On notera encore $D$ les filtrations induites
sur $\scrU\bfv$, $\scrU^{\uphat} \bfv$ et son
dual $\scrO(V)$.
La filtration $n\mapsto D^{-n}\scrO(V)$
de $\scrO(V)$ est croissante, et
$D^1\scrO(V)=0$.

Pour $\bfv$ muni du crochet nul,
$\exp(\bfv,\text{crochet }0)$ est le sch\'ema
pro-vectoriel $\bfv$.

Par A.1, l'application exponentielle $\exp\colon\,
\bfv\to V$ respecte les filtrations $D$ des
alg\`ebres affines.
Par r\'eduction au cas lin\'eaire, on en
d\'eduit que l'alg\`ebre affine de
$V/D^{n+1}V$ est la sous-alg\`ebre de
$\scrO(V)$ engendr\'ee par $D^{-n}\scrO(V)$.
En particulier, $D^0\scrO(V)$ est l'alg\`ebre
affine de $V/D^1V$.

Si $M$ est une repr\'esentation de $V$, les
conditions suivantes sur une filtration $F$ de $M$
sont \'equivalentes:

\medskip\noindent
\underbar{$\,\,\,$}
compatiblit\'e \`a $D$ sur $\bfv$:
$D^p\bfv\otimes F^qM\to F^{p+q}M$,

\noindent
\underbar{$\,\,\,$}
compatibilit\'e \`a $D$ sur $\scrU\bfv$:
$D^p\scrU\bfv\otimes F^qM\to F^{p+q}M$,

\noindent
\underbar{$\,\,\,$}
compatibilit\'e \`a $D$ sur $\scrO(V)$:
$F^pM\to\sum D^{-r}\scrO(V)\otimes F^{p+r}M$

\medskip
La filtration $D$ de $\scrO(V)$ est stable par
translations \`a gauche et \`a droite.
Si $P$ est un $V$-torseur, elle fournit encore
une filtration $D$ sur l'alg\`ebre affine
$\scrO(P)$ de $P$.

\subhead
A.11 Fonctorialit\'es
\endsubhead
(i)\enspace
Soient $\bfv_1$ une sous-alg\`ebre de Lie de
$\bfv$ (resp. un quotient) et $V_1$ le
sous-groupe (resp. quotient) correspondant de
$V$.
Notons $D_1$ la filtration de $\bfv_1$ induite
par (resp. quotient de) la filtration $D$ de
$\bfv$, et les filtrations qui s'en
d\'eduisent sur $\Sym^*\bfv_1$, $\scrU\bfv_1$,
$\scrU^{\uphat}\bfv_1$ et $\scrO(V_1)$.
La puissance symm\'etrique $\Sym^n(\bfv_1)$
est un sous-espace (resp. quotient) de
$\Sym^n(\bfv)$, et sa filtration $D_1$ est
induite par (resp. quotient de) la filtration
$D$ de $\Sym^n(\bfv)$.
Il r\'esulte de A.1 que la m\^{e}me assertion
vaut pour $\scrU v_1$ et $\scrU^{\uphat}v_1$
et que dualement, la filtration $D_1$ du
quotient (resp. sous-espace) $\scrO(V_1)$ de
$\scrO(V)$ est quotient de (resp. induite
par) la filtration $D$ de $\scrO(V)$.

\noindent
(ii)\enspace
Supposons que la filtration $D$ de $\bfv$
provienne d'une graduation
$\bfv=\oplusop_{n\Ge 0}\bfv_n$.
Cette graduation se propage \`a
$\Sym^*(\bfv)$, $\scrU\bfv$,
$\scrU^{\uphat}\bfv$ et $\scrO(V)$.
Pour $\scrU^{\uphat}\bfv$, ``graduation'' et
\`a prendre au sens de ``d\'ecomposition en
produit'' plut\^{o}t que ``d\'ecomposition en
somme''.
La filtration $D$ de $\Sym^n(\bfv)$ provient
de la graduation de $\Sym^n(\bfv)$.
Par A.1, la m\^{e}me assertion vaut pour
$\scrU\bfv$, $\scrU^{\uphat}\bfv$ et
$\scrO(V)$.

\noindent
(iii)\enspace
Par passage \`a la limite, A.10 et A.11
restent valable dans le cas pro-unipotent.

\proclaim{A.12 Exemples} \rm
(i)\enspace
Prenons l'alg\`ebre de Lie pro-unipotente
$\Lib^{\uphat}$ de A.9, et la filtration centrale
descendante $\grZ$.
Elle provient de la graduation pour laquelle
les $x_i$ sont de degr\'e $1$.
Chaque mon\^{o}me $m$ en les $x_i$ d\'efinit
sur $\exp(\Lib^{\uphat})\Subset k\left<\!\left<
(x_i\right>\!\right>$ une fonction $c_m$
``coefficient de $m$ dans $g$'', et
l'alg\`ebre affine a pour base les $c_m$.
Le produit est le produit de m\'elange.
Le degr\'e de $c_m$ est l'oppos\'e de celui de
$m$.

\noindent
(ii)\enspace
Parmi les g\'en\'erateurs $x_i$, distinguons en
un, not\'e $x_0$, et posons $I=\{0\}\amalg J$.
On a 
$$
\Lib:=\Lib((x_i)_{i\in I})=k.x_0\ltimes
\Lib((\ad\,x_0)^n(x_j),\,\,
(n\Ge 0,\,j\in J)).
$$
Soit $D$ la filtration de Lib pour laquelle
$D^p$ est, pour $p\Ge 1$, la s\'erie centrale
descendante du second facteur.
La filtration $D$ provient de la graduation
pour laquelle $x_0$ est de degr\'e $0$ et les
$x_j$ $(j\in J)$ de degr\'e $1$.
Compl\'etant cette filtration, on obtient une
filtration, encore not\'ee $D$, de
$\Lib^{\uphat}$.
La filtration correspondante de l'alg\`ebre
affine
$$
\scrO(\exp(\Lib^{\uphat}))
$$
est d\'eduite de la graduation pour laquelle
le degr\'e de $c_m$ est l'oppos\'e du degr\'e
du mon\^{o}me $m$ en les $x_j$ $(j\in J)$.
\endproclaim

\subhead
A.13
\endsubhead
Soit $G$ un groupe alg\'ebrique produit
semi-direct du groupe multiplicatif $\dbG_m$
par un groupe alg\'ebrique unipotent $V$,
d'alg\`ebre de Lie $\bfv$.
L'action de $\dbG_m$ sur $\bfv$ gradue cette
alg\`ebre de Lie; sur $\bfv_n$,
$\lambda\in\dbG_m$ agit par multiplication par
$\lambda^n$.

Supposons $\bfv$ \`a degr\'es $>0$.
Le foncteur $\rho\mapsto d\rho$ est alors une
\'equivalence de la cat\'egorie des
repr\'esentations lin\'eaires de $G$ avec
celle des espaces vectoriels gradu\'es $E$,
muni d'une action compatible aux graduations
de l'alg\`ebre de Lie $\bfv$.
La graduation de $E$ donne l'action de
$\dbG_m$ et l'action de $\bfv$,
automatiquement nilpotente, donne celle de $V$.

Notons $k(a)$ l'espace vectoriel $k$
plac\'e en degr\'e $a$.
Pour $a<b$, une extension de $k(a)$ par
$k(b)$ dans la cat\'egorie des
repr\'esentations de $G$ est donn\'ee par un
homomorphisme $\bfv\to
\Hom(k(a),k(b))$, i.e.
par un \'el\'ement du dual de la partie
homog\`ene de degre\'e $b-a$ de $\bfv^{\ab}:=
\bfv/[\bfv,\bfv]$:
$$
\Ext_{\Rep(G)}^1(k(a),k(b))=
(\bfv^{\ab})_{b-a}{\upvee}.
\tag{A.13.1}
$$

\subhead
A.14
\endsubhead
Soit $G$ un sch\'ema en groupes produit
semi-direct de  $\dbG_m$ par $V$
pro-unipotent.
Ecrivons-le comme limite projective de
quotients par des sous-groupes distingu\'es
contenus dans $V$:
$$
G=\dbG_m\ltimes V=\lim\,\dbG_m\ltimes
V_\alpha.
$$

Supposons les alg\`ebres de Lie gradu\'ees
$\bfv_\alpha$ des $V_\alpha$ (A.13) \`a
degr\'e $>0$.
Passant \`a la limite, on d\'eduit de (A.13.1)
que
$$
\Ext_{\Rep(G)}^1(k(a),k(b))=
\colim(v_\alpha^{\ab})_{b-a}{\upvee}.
\tag{A.14.1}
$$

Supposons en outre que pour tout,
$$
\dim\,\Ext^1(k(0),k(n))<\infty.
$$
Les limites projectives (A.14.1) sont alors
stationnaires et, puisque si on rel\`eve dans les
$(\bfv_\alpha)_n$ des bases des
$(\bfv_\alpha^{\ab})_n$, on obtient un syst\`eme
g\'en\'erateur de $\bfv_\alpha$, la dimension de
$(\bfv_\alpha)_n$ est, pour chaque $n$, born\'ee
ind\'ependamment de $n$.
On notera $\bfv_{\GR}$ l'alg\`ebre de Lie
gradu\'ee de composantes homog\`enes les
limites projectives (stationnaires) des
composantes homog\`enes des $v_\alpha$.
Passant \`a la limite dans A.13, on obtient:

\proclaim{Proposition A.15}
Sous les hypoth\`eses de A.14:
$G=\dbG_m\ltimes V$, Lie $V$ \`a degr\'es
$>0$, $\Ext_{\Rep(G)}^1(k(0),k(n))$
de dimension finie, le foncteur $\rho\mapsto
d\rho$ est une \'equivalence de la cat\'egorie
des repr\'esentations de $G$ avec celle des
espaces vectoriels gradu\'es munis d'une 
action de l'alg\`ebre de Lie gradu\'ee
$\bfv_{\GR}$.
On a
$$
\Ext_{\Rep(G)}^1(k(0),k(n))=
(\bfv_{\GR}^{\ab})_n\upvee.
$$
\endproclaim

